% !TeX encoding = UTF-8
\documentclass[variant=guia,cover=institutional,mode=screen]{ua-engineering-v6-article}
% !TeX encoding = UTF-8
\UASetUniversity{Universidad Austral}
\UASetFaculty{Facultad de Ingeniería}
\UASetDepartment{Departamento de Computación}
\UASetCampus{Pilar}
\UASetCareer{Ingeniería Informática}
\UASetCommission{Comisión A}
\UASetProfessor{Equipo docente}
\UASetSemester{Segundo cuatrimestre 2026}
\UASetCourse{Sistemas Distribuidos}
\UASetDate{2026}


\UASetClassNumber{14}
\UASetClassTitle{Incidentes y cultura de confiabilidad}
\UASetDocKind{Guía de ejercicios}

\title{Clase 14 --- Guía de ejercicios}
\author{\UAProfessor}
\date{\UADate}

\begin{document}

\section{Indicaciones generales}
Resuelva cada ejercicio declarando supuestos, modelo de fallas, garantía y trade-off. Cuando corresponda, construya una ejecución verificable y vincule la respuesta con evidencia observable.

\section{Nivel 1: fundamentos y reconocimiento}
\begin{uaexercise}
Defina incidente desde la perspectiva de confiabilidad operativa y explique por qué no debe confundirse con cualquier bug menor.
\end{uaexercise}

\begin{uaexercise}
Explique por qué en sistemas complejos los incidentes deben considerarse eventos esperables y no rarezas imposibles.
\end{uaexercise}

\begin{uaexercise}
Distinga claramente entre respuesta al incidente y aprendizaje posterior al incidente.
\end{uaexercise}

\begin{uaexercise}
Explique por qué una organización puede tener buenos sistemas técnicos y, aun así, mala gestión de incidentes.
\end{uaexercise}

\begin{uaexercise}
Describa las fases principales de la gestión de incidentes.
\end{uaexercise}


\section{Nivel 2: ejecuciones y fallas}
\begin{uaexercise}
Explique qué criterios deberían usarse para priorizar un incidente.
\end{uaexercise}

\begin{uaexercise}
Explique qué ventajas tiene explicitar roles durante un incidente.
\end{uaexercise}

\begin{uaexercise}
Construya un ejemplo donde la falta de coordinación empeore significativamente el incidente.
\end{uaexercise}

\begin{uaexercise}
Defina postmortem y explique cuál es su propósito.
\end{uaexercise}

\begin{uaexercise}
Explique qué información mínima debe contener un buen postmortem.
\end{uaexercise}


\section{Nivel 3: diseño y comparación}
\begin{uaexercise}
Explique por qué el timeline es una pieza central del postmortem.
\end{uaexercise}

\begin{uaexercise}
Construya un ejemplo donde un postmortem superficial impida aprendizaje real.
\end{uaexercise}

\begin{uaexercise}
Explique qué significa blameless postmortem.
\end{uaexercise}

\begin{uaexercise}
Explique por qué blameless no significa ausencia de responsabilidad.
\end{uaexercise}

\begin{uaexercise}
Explique por qué la idea de una única causa raíz puede ser engañosa en sistemas distribuidos.
\end{uaexercise}


\section{Nivel 4: integración y defensa}
\begin{uaexercise}
Construya una cadena causal realista para un incidente de latencia en cascada.
\end{uaexercise}

\begin{uaexercise}
Explique qué hace buena a una acción correctiva posterior a un incidente.
\end{uaexercise}

\begin{uaexercise}
Dé tres ejemplos de malas acciones de postmortem y justifique por qué lo son.
\end{uaexercise}

\begin{uaexercise}
Explique cómo distinguir entre fix local y aprendizaje organizacional.
\end{uaexercise}

\begin{uaexercise}
Explique por qué cerrar un incidente sin cerrar el seguimiento del postmortem es una mala práctica.
\end{uaexercise}


\end{document}
